\chapter{Reference Card}\label{app:cheatsheet}

\section{Modes}

\begin{keyslong}[0.24]
	\key{<Esc>}   & Return to Normal mode from anywhere        \\
	\key{i} \key{a} \key{I} \key{A} \key{o} \key{O} & Enter Insert mode \\
	\key{v} \key{V} \key{<C-v>} & Charwise / linewise / blockwise Visual \\
	\key{:}       & Command-line mode                          \\
	\key{R}       & Replace mode                               \\
	\key{<C-\textbackslash>} \key{<C-n>} & Leave Terminal mode \\
\end{keyslong}

\section{Motions}

\begin{keyslong}[0.24]
	\key{h} \key{j} \key{k} \key{l}   & Left, down, up, right                  \\
	\key{w} \key{b} \key{e}           & Word forward, back, end                \\
	\key{W} \key{B} \key{E}           & Same, by whitespace-delimited WORD     \\
	\key{0} \key{\textasciicircum{}} \key{\$}         & Line start, first non-blank, line end  \\
	\key{f} \key{x} / \key{t} \key{x} & To / till the next \texttt{x}; \key{;} and \key{,} repeat \\
	\key{g} \key{g} / \key{G}         & Top / bottom of the file               \\
	\key{42} \key{G}                  & Line 42                                \\
	\key{\{} \key{\}}                 & Previous / next paragraph              \\
	\key{\%}                          & Matching bracket                       \\
	\key{H} \key{M} \key{L}           & Top / middle / bottom of the screen    \\
	\key{<C-d>} \key{<C-u>}           & Half page down / up                    \\
	\key{z} \key{z}                   & Centre the cursor line                 \\
	\key{<C-o>} \key{<C-i>}           & Back / forward in the jump list        \\
	\key{m} \key{a} then \key{`} \key{a} & Set and return to mark \texttt{a}   \\
\end{keyslong}

\section{Operators and Text Objects}

\begin{keyslong}[0.24]
	\key{d} \key{c} \key{y}           & Delete, change, yank                   \\
	\key{>} \key{<} \key{=}           & Indent, dedent, re-indent              \\
	\key{g} \key{u} \key{g} \key{U} \key{g} \key{\textasciitilde} & Lower, upper, toggle case \\
	\key{g} \key{q}                   & Reformat to \cmd{'textwidth'}          \\
	\key{d} \key{d} \key{y} \key{y} \key{c} \key{c} & Operate on the whole line \\
	\key{D} \key{C} \key{Y}           & Operate to the end of the line         \\
	\key{x} / \key{X}                 & Delete the character under / before the cursor \\
	\key{r} \key{x}                   & Replace one character with \texttt{x}  \\
	\key{J}                           & Join with the next line                \\
	\key{i} \key{w} / \key{a} \key{w} & Inner / a word                         \\
	\key{i} \key{p} / \key{a} \key{p} & Inner / a paragraph                    \\
	\key{i} \key{"} \key{i} \key{(} \key{i} \key{\{} \key{i} \key{[} \key{i} \key{t} & Inside quotes, parens, braces, brackets, tags \\
	\key{a} \key{"} \key{a} \key{(} \dots & The same, including the delimiters \\
	\key{.}                           & Repeat the last change                 \\
	\key{u} / \key{<C-r>}             & Undo / redo                            \\
\end{keyslong}

\section{Registers, Macros, Search}

\begin{keyslong}[0.24]
	\key{"} \key{a} \key{y} \key{y}   & Yank a line into register \texttt{a}   \\
	\key{"} \key{a} \key{p}           & Paste register \texttt{a}              \\
	\key{"} \key{0} \key{p}           & Paste the last \emph{yank}, not the last delete \\
	\key{"} \key{+} \key{y} / \key{"} \key{+} \key{p} & System clipboard      \\
	\key{"} \key{\_} \key{d}          & Delete without touching any register   \\
	\cmd{:registers}                  & Show every register                    \\
	\key{q} \key{a} \dots \key{q}     & Record a macro into register \texttt{a} \\
	\key{@} \key{a} / \key{@} \key{@} & Play it / replay the last one          \\
	\key{/} \texttt{pat} \key{<CR>}   & Search forward; \key{n} and \key{N} repeat \\
	\key{*} / \key{\#}                & Search for the word under the cursor   \\
	\cmd{:noh}                        & Clear search highlighting              \\
	\cmd{:\%s/old/new/gc}             & Substitute in the file, confirming each \\
	\cmd{:g/pat/normal @q}            & Run macro \texttt{q} on matching lines \\
\end{keyslong}

\section{Files, Buffers, Windows}

\begin{keyslong}[0.24]
	\cmd{:w} \cmd{:q} \cmd{:wq} \cmd{:q!} & Write, quit, both, force-quit      \\
	\cmd{:e file} / \cmd{:ls} / \cmd{:b N} & Open, list, switch buffer         \\
	\key{<C-\textasciicircum{}>}                      & Toggle with the alternate buffer       \\
	\key{<C-w>} \key{s} / \key{<C-w>} \key{v} & Split horizontally / vertically \\
	\key{<C-w>} \key{h} \key{j} \key{k} \key{l} & Move between windows         \\
	\key{<C-w>} \key{c} / \key{<C-w>} \key{o} & Close this / all other windows \\
	\cmd{:tabnew} / \key{g} \key{t}   & New tab / next tab                     \\
	\cmd{:copen} / \cmd{:cn} / \cmd{:cp} & Quickfix list: open, next, previous \\
	\cmd{:\%!cmd}                     & Filter the buffer through a shell command \\
	\cmd{:help <topic>}               & The manual                             \\
\end{keyslong}

\section{Neovim}

\begin{keyslong}[0.30]
	\cmd{:checkhealth}                & Diagnose the whole installation        \\
	\cmd{:Lazy}                       & Plugin manager dashboard               \\
	\cmd{:Mason}                      & Install language servers and formatters \\
	\cmd{:LspInfo}                    & Which servers serve this buffer        \\
	\cmd{:Inspect} / \cmd{:InspectTree} & Highlight groups / Treesitter tree   \\
	\cmd{:terminal}                   & Shell in a buffer                      \\
	\cmd{:lua} \dots{} / \cmd{:=expr} & Execute / print Lua                    \\
	\cmd{:verbose set opt?}           & Which file set this option             \\
	\cmd{:source \%}                  & Re-run the Lua file being edited       \\
	\key{K} \key{g} \key{d} \key{g} \key{r} & LSP hover, definition, references \\
	\key{]} \key{d} / \key{[} \key{d} & Next / previous diagnostic             \\
\end{keyslong}

\chapter{A Minimal \texttt{init.lua}}\label{app:init}

A self-contained starting point, for when a distribution is more than you want.
Put it in \file{\textasciitilde/.config/nvim/init.lua}; it bootstraps its own
plugin manager on first launch.

\begin{luacode}
-- ── Leaders must be set before any plugin loads ──────────────────────────
vim.g.mapleader = " "
vim.g.maplocalleader = "\\"

-- ── Options ──────────────────────────────────────────────────────────────
local opt = vim.opt
opt.number = true
opt.relativenumber = true
opt.expandtab = true
opt.shiftwidth = 2
opt.tabstop = 2
opt.smartindent = true
opt.ignorecase = true
opt.smartcase = true
opt.undofile = true          -- persistent undo tree
opt.scrolloff = 8
opt.signcolumn = "yes"       -- stop the text jumping when a sign appears
opt.splitright = true
opt.splitbelow = true
opt.termguicolors = true

-- ── Mappings ─────────────────────────────────────────────────────────────
local map = vim.keymap.set
map("n", "<leader>w", "<cmd>write<cr>", { desc = "Save" })
map("n", "<leader>h", "<cmd>nohlsearch<cr>", { desc = "Clear highlight" })
map("n", "<C-h>", "<C-w>h", { desc = "Left window" })
map("n", "<C-j>", "<C-w>j", { desc = "Lower window" })
map("n", "<C-k>", "<C-w>k", { desc = "Upper window" })
map("n", "<C-l>", "<C-w>l", { desc = "Right window" })
map("v", "<", "<gv", { desc = "Dedent, keep selection" })
map("v", ">", ">gv", { desc = "Indent, keep selection" })
map("t", "<Esc><Esc>", "<C-\\><C-n>", { desc = "Leave terminal mode" })

-- ── Autocommands ─────────────────────────────────────────────────────────
vim.api.nvim_create_autocmd("TextYankPost", {
  group = vim.api.nvim_create_augroup("yank_highlight", { clear = true }),
  callback = function() vim.highlight.on_yank({ timeout = 200 }) end,
})

-- ── Plugins (lazy.nvim, bootstrapped on first run) ───────────────────────
local lazypath = vim.fn.stdpath("data") .. "/lazy/lazy.nvim"
if not (vim.uv or vim.loop).fs_stat(lazypath) then
  vim.fn.system({ "git", "clone", "--filter=blob:none", "--branch=stable",
    "https://github.com/folke/lazy.nvim.git", lazypath })
end
vim.opt.rtp:prepend(lazypath)

require("lazy").setup({
  { "folke/tokyonight.nvim", priority = 1000,
    config = function() vim.cmd.colorscheme("tokyonight") end },
  { "nvim-treesitter/nvim-treesitter", build = ":TSUpdate",
    opts = { highlight = { enable = true }, indent = { enable = true } },
    config = function(_, opts)
      require("nvim-treesitter.configs").setup(opts)
    end },
  { "neovim/nvim-lspconfig" },
  { "lervag/vimtex", ft = "tex", init = function()
      vim.g.vimtex_view_method = "zathura"
    end },
})
\end{luacode}
